% =============================================================================
\chapter{Time-Dependent and Nonlinear Problems}
\label{ch:time_nonlinear}
% =============================================================================

This chapter extends finite element discretization beyond stationary linear
boundary-value problems. It develops semidiscrete formulations for transient
PDEs, introduces nonlinear residual equations and consistent tangents, and
surveys the time integration strategies most relevant to multiphysics
simulation after the spectral theory of the preceding chapter.

\section{Semi-discrete finite element systems}
\label{sec:time_nonlinear:semidiscrete}

% Method of lines, mass matrices, and differential-algebraic structure

\section{One-step and multistep time integration}
\label{sec:time_nonlinear:time_integrators}

% Backward Euler, BDF, theta-methods, Newmark, and generalized-alpha

\section{Energy, stability, and numerical dissipation}
\label{sec:time_nonlinear:stability}

% A-stability, L-stability, conservation, and algorithmic damping

\section{Nonlinear residual equations and linearization}
\label{sec:time_nonlinear:linearization}

% Gateaux derivatives, consistent tangents, and incremental formulations

\section{Newton and quasi-Newton methods}
\label{sec:time_nonlinear:newton}

% Convergence criteria, globalization, line search, and inexact solves

\section{Coupled fields and split versus monolithic solves}
\label{sec:time_nonlinear:coupling}

% Partitioned versus monolithic time advancement for multiphysics systems

\section{Adaptive stepping and error control}
\label{sec:time_nonlinear:adaptive}

% Local truncation error estimates, event handling, and multirate ideas
